% ============================================================
\chapter{Fonctions Caract\'eristiques}
\label{ch:caracteristiques}
% ============================================================

\begin{intuition}
La fonction caract\'eristique est un outil analytique puissant qui transforme les
probl\`emes probabilistes en probl\`emes d'analyse complexe. Elle existe toujours
(contrairement \`a la MGF), caract\'erise enti\`erement la loi, et permet de
d\'emontrer \'el\'egamment le th\'eor\`eme central limite.
\end{intuition}

% ------------------------------------------------------------
\section{D\'efinition et premi\`eres propri\'et\'es}
% ------------------------------------------------------------

\begin{definition}[Fonction caract\'eristique]
La \textbf{fonction caract\'eristique} d'une variable al\'eatoire r\'eelle $X$ est~:
\[
    \varphi_X(t) = \E[e^{itX}] = \E[\cos(tX)] + i\,\E[\sin(tX)], \quad t \in \R.
\]
\end{definition}

\begin{theoreme}[Propri\'et\'es fondamentales]
\label{thm:prop_fc}
Soit $\varphi_X$ la fonction caract\'eristique de $X$. Alors~:
\begin{enumerate}[label=(\roman*)]
    \item $\varphi_X(0) = 1$.
    \item $\abs{\varphi_X(t)} \leq 1$ pour tout $t \in \R$.
    \item $\varphi_X$ est \textbf{uniform\'ement continue} sur $\R$.
    \item $\varphi_X(-t) = \overline{\varphi_X(t)}$ (conjugu\'e complexe).
    \item Si $Y = aX + b$, alors $\varphi_Y(t) = e^{ibt}\, \varphi_X(at)$.
    \item Si $X \indep Y$, alors $\varphi_{X+Y}(t) = \varphi_X(t)\, \varphi_Y(t)$.
\end{enumerate}
\end{theoreme}

\begin{proof}[Preuve de (iii)]
\begin{align*}
    \abs{\varphi_X(t+h) - \varphi_X(t)}
    &= \abs{\E[e^{itX}(e^{ihX} - 1)]}
    \leq \E[\abs{e^{ihX} - 1}].
\end{align*}
Comme $\abs{e^{ihX} - 1} \leq 2$ et $e^{ihX} - 1 \to 0$ quand $h \to 0$,
la convergence domin\'ee donne $\E[\abs{e^{ihX}-1}] \to 0$ uniform\'ement en $t$.
\end{proof}

% ------------------------------------------------------------
\section{Calcul des moments}
% ------------------------------------------------------------

\begin{theoreme}[Moments et d\'eriv\'ees]
Si $\E[\abs{X}^n] < \infty$, alors $\varphi_X$ est $n$ fois d\'erivable et~:
\[
    \varphi_X^{(k)}(0) = i^k\, \E[X^k], \quad k = 0, 1, \ldots, n.
\]
En particulier~:
$\E[X] = \varphi_X'(0)/i$ et
$\E[X^2] = -\varphi_X''(0)$, d'o\`u
$\Var(X) = -\varphi_X''(0) - [\varphi_X'(0)/i]^2$.
\end{theoreme}

% ------------------------------------------------------------
\section{Fonctions caract\'eristiques des lois classiques}
% ------------------------------------------------------------

\begin{formules}[title=\textbf{Table des fonctions caract\'eristiques}]
\begin{center}
\renewcommand{\arraystretch}{1.5}
\begin{tabular}{lc}
    \toprule
    \textbf{Loi de $X$} & $\varphi_X(t)$ \\
    \midrule
    $\text{Bernoulli}(p)$ & $1 - p + pe^{it}$ \\
    $\text{Bin}(n,p)$ & $(1-p+pe^{it})^n$ \\
    $\text{Poisson}(\lambda)$ & $\exp(\lambda(e^{it}-1))$ \\
    $\text{Geom}(p)$ & $\dfrac{pe^{it}}{1-(1-p)e^{it}}$ \\[4pt]
    $\mathcal{U}([a,b])$ & $\dfrac{e^{itb}-e^{ita}}{it(b-a)}$ \\[4pt]
    $\text{Exp}(\lambda)$ & $\dfrac{\lambda}{\lambda - it}$ \\[4pt]
    $\mathcal{N}(\mu, \sigma^2)$ & $\exp\!\left(i\mu t - \dfrac{\sigma^2 t^2}{2}\right)$ \\[4pt]
    $\Gamma(\alpha, \lambda)$ & $\left(\dfrac{\lambda}{\lambda - it}\right)^\alpha$ \\[4pt]
    $\text{Cauchy}(0,1)$ & $e^{-\abs{t}}$ \\
    \bottomrule
\end{tabular}
\end{center}
\end{formules}

\begin{exemple}[Calcul pour la loi normale]
Si $X \sim \mathcal{N}(0,1)$~:
\begin{align*}
    \varphi_X(t) &= \int_{-\infty}^{+\infty} e^{itx}\, \frac{1}{\sqrt{2\pi}}\, e^{-x^2/2}\, dx
    = \frac{1}{\sqrt{2\pi}} \int_{-\infty}^{+\infty} e^{-(x^2 - 2itx)/2}\, dx \\
    &= \frac{1}{\sqrt{2\pi}} \int_{-\infty}^{+\infty}
    e^{-(x-it)^2/2}\, e^{-t^2/2}\, dx
    = e^{-t^2/2}.
\end{align*}
Pour $X \sim \mathcal{N}(\mu, \sigma^2)$, $X = \sigma Z + \mu$, donc
$\varphi_X(t) = e^{i\mu t}\, e^{-\sigma^2 t^2/2} = e^{i\mu t - \sigma^2 t^2/2}$.
\end{exemple}

% ------------------------------------------------------------
\section{Th\'eor\`emes d'unicit\'e et d'inversion}
% ------------------------------------------------------------

\begin{theoreme}[Th\'eor\`eme d'unicit\'e]
\label{thm:unicite}
Deux variables al\'eatoires $X$ et $Y$ ont m\^eme loi si et seulement si
$\varphi_X(t) = \varphi_Y(t)$ pour tout $t \in \R$.
\end{theoreme}

\begin{theoreme}[Formule d'inversion]
\label{thm:inversion}
Si $\varphi_X \in L^1(\R)$ (i.e.\ $\int \abs{\varphi_X(t)}\, dt < \infty$),
alors $X$ admet une densit\'e continue donn\'ee par~:
\[
    f_X(x) = \frac{1}{2\pi} \int_{-\infty}^{+\infty} e^{-itx}\, \varphi_X(t)\, dt.
\]
\end{theoreme}

\begin{remarque}
La formule d'inversion est l'analogue probabiliste de la transform\'ee de Fourier
inverse. La fonction caract\'eristique n'est autre que la transform\'ee de Fourier
de la loi de $X$.
\end{remarque}

% ------------------------------------------------------------
\section{Th\'eor\`eme de continuit\'e de L\'evy}
% ------------------------------------------------------------

\begin{theoreme}[Th\'eor\`eme de continuit\'e de L\'evy]
\label{thm:levy}
Soit $(X_n)$ une suite de variables al\'eatoires et $\varphi_n = \varphi_{X_n}$.
\begin{enumerate}[label=(\roman*)]
    \item Si $X_n \xrightarrow{\mathcal{L}} X$, alors $\varphi_n(t) \to \varphi_X(t)$
    pour tout $t$.
    \item R\'eciproquement, si $\varphi_n(t) \to \psi(t)$ pour tout $t$ et si $\psi$
    est \textbf{continue en 0}, alors $\psi$ est la fonction caract\'eristique d'une
    variable $X$ et $X_n \xrightarrow{\mathcal{L}} X$.
\end{enumerate}
\end{theoreme}

\begin{remarque}
C'est ce th\'eor\`eme qui est utilis\'e dans la preuve du TCL (chapitre~10)~: on
montre la convergence ponctuelle de $\varphi_{Z_n}(t) \to e^{-t^2/2}$, et la
continuit\'e en 0 de la limite est \'evidente.
\end{remarque}

% ------------------------------------------------------------
\section{Application : preuve de stabilit\'e de la loi normale}
% ------------------------------------------------------------

\begin{theoreme}[Stabilit\'e de la loi normale]
Si $X_1 \sim \mathcal{N}(\mu_1, \sigma_1^2)$ et
$X_2 \sim \mathcal{N}(\mu_2, \sigma_2^2)$ sont ind\'ependantes, alors~:
\[
    X_1 + X_2 \sim \mathcal{N}(\mu_1 + \mu_2, \sigma_1^2 + \sigma_2^2).
\]
\end{theoreme}

\begin{proof}
\begin{align*}
    \varphi_{X_1+X_2}(t) &= \varphi_{X_1}(t)\, \varphi_{X_2}(t) \\
    &= e^{i\mu_1 t - \sigma_1^2 t^2/2}\, e^{i\mu_2 t - \sigma_2^2 t^2/2} \\
    &= e^{i(\mu_1+\mu_2)t - (\sigma_1^2+\sigma_2^2)t^2/2},
\end{align*}
qui est la FC de $\mathcal{N}(\mu_1+\mu_2, \sigma_1^2+\sigma_2^2)$. Le r\'esultat
suit par unicit\'e.
\end{proof}

% ------------------------------------------------------------
\section{Application : stabilit\'e de la loi de Poisson}
% ------------------------------------------------------------

\begin{proposition}
Si $X_1 \sim \text{Poisson}(\lambda_1)$ et $X_2 \sim \text{Poisson}(\lambda_2)$
sont ind\'ependantes, alors
$X_1 + X_2 \sim \text{Poisson}(\lambda_1 + \lambda_2)$.
\end{proposition}

\begin{proof}
$\varphi_{X_1+X_2}(t) = e^{\lambda_1(e^{it}-1)} \cdot e^{\lambda_2(e^{it}-1)}
= e^{(\lambda_1+\lambda_2)(e^{it}-1)}$, qui est la FC de
$\text{Poisson}(\lambda_1+\lambda_2)$.
\end{proof}

% ------------------------------------------------------------
\section{Fonction caract\'eristique et convergence en loi}
% ------------------------------------------------------------

\begin{exemple}[Convergence de Poisson vers la normale]
Soit $X_n \sim \text{Poisson}(n)$. Posons $Z_n = (X_n - n)/\sqrt{n}$.
\begin{align*}
    \varphi_{Z_n}(t) &= e^{-itn/\sqrt{n}}\, \varphi_{X_n}(t/\sqrt{n})
    = e^{-it\sqrt{n}}\, e^{n(e^{it/\sqrt{n}}-1)} \\
    &= \exp\!\left(n\!\left(e^{it/\sqrt{n}} - 1 - \frac{it}{\sqrt{n}}\right)\right).
\end{align*}
Par d\'eveloppement~: $e^{it/\sqrt{n}} = 1 + it/\sqrt{n} - t^2/(2n) + o(1/n)$, d'o\`u~:
\[
    \varphi_{Z_n}(t) = \exp\!\left(n\!\left(-\frac{t^2}{2n} + o(1/n)\right)\right)
    \to e^{-t^2/2}.
\]
Donc $Z_n \xrightarrow{\mathcal{L}} \mathcal{N}(0,1)$ par le th\'eor\`eme de L\'evy.
\end{exemple}

% ------------------------------------------------------------
\section{Exercices}
% ------------------------------------------------------------

\begin{exercice}
Calculer la fonction caract\'eristique de la loi $\mathcal{U}(\{-1, 0, 1\})$.
\end{exercice}

\begin{exercice}
Soit $X \sim \text{Exp}(\lambda)$. Calculer $\varphi_X(t)$ et en d\'eduire
$\E[X]$, $\E[X^2]$, $\Var(X)$.
\end{exercice}

\begin{exercice}
Montrer que $\varphi_X(t) = e^{-\abs{t}}$ est la fonction caract\'eristique de la
loi de Cauchy standard. En d\'eduire que si $X_1, X_2$ sont i.i.d.\ Cauchy standard,
alors $(X_1+X_2)/2$ est encore Cauchy standard.
\end{exercice}

\begin{exercice}
En utilisant la formule d'inversion, retrouver la densit\'e de la loi $\mathcal{N}(0,1)$
\`a partir de sa FC $\varphi(t) = e^{-t^2/2}$.
\end{exercice}

\begin{exercice}
Montrer que la FC de la loi $\Gamma(\alpha, \lambda)$ est
$(\lambda/(\lambda - it))^\alpha$. En d\'eduire que la somme de $n$ exponentielles
i.i.d.\ de param\`etre $\lambda$ suit une loi $\Gamma(n, \lambda)$.
\end{exercice}

\begin{exercice}
Soit $X_n \sim \text{Bin}(n, \lambda/n)$ avec $\lambda > 0$ fix\'e. En calculant
$\varphi_{X_n}(t)$, retrouver l'approximation de Poisson~:
$X_n \xrightarrow{\mathcal{L}} \text{Poisson}(\lambda)$.
\end{exercice}

\begin{exercice}
Montrer que si $\varphi_X$ est \`a valeurs r\'eelles, alors $X$ et $-X$ ont m\^eme
loi (la loi de $X$ est sym\'etrique autour de 0).
\end{exercice}

\begin{exercice}
Soit $(X_n)$ i.i.d.\ avec $\E[X_1] = 0$, $\Var(X_1) = \sigma^2$ et
$\varphi_{X_1}(t) = 1/(1+t^2)$ (Cauchy). Montrer que $\bar{X}_n$ a m\^eme loi
que $X_1$ et que le TCL ne s'applique pas (pourquoi~?).
\end{exercice}
